\begin{thebibliography}{999}
\addcontentsline{toc}{chapter}{参考文献}

\bibitem{bib1} Adámek, J., and Rosický, J. 1994. \emph{Locally Presentable and
  Accessible Categories}. Cambridge University Press.

\bibitem{bib2} Anel, M., and Lejay, D. 2018. Exponentiable higher toposes.
  arXiv:1802.10425.

\bibitem{bib3} Antolin-Camarena, O. 2016. A whirlwind tour of the world of
  $(\infty,1)$-categories. Page 15--62 of: Bárcenas, N., Galaz-Garcia, F., and
  Rocha, M. M. (eds), \emph{Mexican Mathematicians Abroad: Recent
  Contributions}. Contemporary Mathematics, vol.~657. American Mathematical
  Society.

\bibitem{bib4} Ara, D. 2014. Higher quasi-categories vs higher Rezk spaces.
  \emph{J. K-Theory}, \textbf{14}(3), 701--749.

\bibitem{bib5} Ayala, D., and Francis, J. 2020. Fibrations of
  $\infty$-categories. \emph{High. Struct.}, \textbf{4}(1), 168--265.

\bibitem{bib6} Barwick, C. 2005. \emph{($\infty,n$)-Cat as a closed model
  category}. Ph.D. thesis, University of Pennsylvania.

\bibitem{bib7} Barwick, C., Glasman, S., and Nardin, D. 2018. Dualizing
  cartesian and cocartesian fibrations. \emph{Theory Appl. Categ.},
  \textbf{33}(4), 67--94.

\bibitem{bib8} Beke, T. 2001. Sheafifiable homotopy model categories. \emph{Math.
  Proc. Cambridge Philos. Soc.}, \textbf{129}(03), 447--475.

\bibitem{bib9} Bénabou, J. 1967. Introduction to Bicategories. Page 1--77 of:
  \emph{Reports of the Midwest Category Seminar}. Lecture Notes in Math.,
  no.~47. Springer-Verlag.

\bibitem{bib10} Bénabou, J. 1973. \emph{Les distributeurs}. Tech. rep.
  Université Catholique de Louvain, Institut de Mathématique Pure et
  Appliquée.

\bibitem{bib11} Berger, C. 2007. Iterated wreath product of the simplex category
  and iterated loop spaces. \emph{Adv. Math.}, \textbf{213}(1), 230--270.

\bibitem{bib12} Berger, C., and Moerdijk, I. 2008. On an extension of the notion
  of Reedy category. \emph{Mathe. Z.}, \textbf{269}, 1--28.

\bibitem{bib13} Bergner, J. E. 2007a. A characterization of fibrant Segal
  categories. \emph{Proc. Amer. Math. Soc.}, \textbf{135}(12), 4031--4037.

\bibitem{bib14} Bergner, J. E. 2007b. Three models for the homotopy theory of
  homotopy theories. \emph{Topology}, \textbf{46}(4), 397--436.

\bibitem{bib15} Bergner, J. E. 2018. \emph{An Introduction to
  ($\infty$,1)-Categories}. Cambridge University Press.

\bibitem{bib16} Bergner, J. E., and Rezk, C. 2013. Reedy categories and the
  $\Theta$-construction. \emph{Mathe. Z.}, \textbf{274}(1--2), 499--514.

\bibitem{bib17} Betti, R., Carboni, A., Street, R., and Walters, R. 1983.
  Variation through enrichment. \emph{J. Pure Appl. Algebra}, \textbf{29},
  109--127.

\bibitem{bib18} Bird, G. J., Kelly, G. M., Power, A. J., and Street, R. H. 1989.
  Flexible Limits for 2-Categories. \emph{J. Pure Appl. Algebra},
  \textbf{61}(1), 1--27.

\bibitem{bib19} Blanc, G. 1978. Équivalence naturelle et formules logiques en
  théorie des catégories. \emph{Archiv math. Logik}, \textbf{19}, 131--137.

\bibitem{bib20} Blumberg, A., and Mandell, M. 2011. Algebraic K-theory and
  abstract homotopy theory. \emph{Adv. Math.}, \textbf{226}, 3760--3812.

\bibitem{bib21} Boardman, J. M., and Vogt, R. 1973. \emph{Homotopy Invariant
  Algebraic Structures on Topological Spaces}. Lecture Notes in Mathematics,
  vol.~347. Springer-Verlag.

\bibitem{bib22} Bourke, J., Lack, S., and Vokřínek, L. 2020. Adjoint functor
  theorems for homotopically enriched categories. arXiv:2006.07843.

\bibitem{bib23} Brown, K. S. 1973. Abstract homotopy theory and generalized sheaf
  cohomology. \emph{Trans. Amer. Math. Soc.}, \textbf{186}, 419--458.

\bibitem{bib24} Burroni, A. 1971. T-catégories (catégories dans un triple).
  \emph{Cah. Topol. Géom. Différ. Catég.}, \textbf{12}, 215--321.

\bibitem{bib25} Cartmell, J. 1986. Generalised algebraic theories and contextual
  categories. \emph{Ann. Pure Appl. Logic}, \textbf{32}, 209--243.

\bibitem{bib26} Cheng, E., Gurski, N., and Riehl, E. 2014. Cyclic
  multicategories, multivariable adjunctions and mates. \emph{J. K-Theory},
  \textbf{13}, 337--396.

\bibitem{bib27} Cisinski, D. C. 2006. \emph{Les préfaisceaux comme modèles des
  types d'homotopie}. Astérisque, vol.~308. Soc. Math. France.

\bibitem{bib28} Cisinski, D. C. 2019. \emph{Higher Categories and Homotopical
  Algebra}. Cambridge Studies in Advanced Mathematics, vol.~180. Cambridge
  University Press.

\bibitem{bib29} Cordier, J. M. 1982. Sur la notion de diagramme homotopiquement
  cohérent. \emph{Cah. Topol. Géom. Différ. Catég.}, \textbf{23}, 93--112.

\bibitem{bib30} Cordier, J. M., and Porter, T. 1986. Vogt's theorem on categories
  of homotopy coherent diagrams. \emph{Math. Proc. Cambridge Philos. Soc.},
  \textbf{100}, 65--90.

\bibitem{bib31} Cruttwell, G. S. H. 2008. \emph{Normed Spaces and the Change of
  Base for Enriched Categories}. Ph.D. thesis, Dalhousie University.

\bibitem{bib32} Cruttwell, G. S. H., and Shulman, M. A. 2010. A unified framework
  for generalized multicategories. \emph{Theory Appl. Categ.},
  \textbf{24}(21), 580--655.

\bibitem{bib33} Day, B. 1970. On closed categories of functors. Page 1--38 of:
  \emph{Reports of the Midwest Category Seminar IV}. Lecture Notes in
  Mathematics, vol.~137. Springer-Verlag.

\bibitem{bib34} de Brito, P. B. 2018. Segal objects and the Grothendieck
  construction. Page 19--44 of: Ausoni, C., Hess, K., Johnson, B., Moerdijk,
  I., and Scherer, J. (eds), \emph{An Alpine Bouquet of Algebraic Topology}.
  Contemporary Mathematics, vol.~708. American Mathematical Society.

\bibitem{bib35} Dugger, D. 2001. Replacing model categories with simplicial ones.
  \emph{Trans. Amer. Math. Soc.}, \textbf{353}(12), 5003--5028.

\bibitem{bib36} Dugger, D., and Spivak, D. I. 2011. Mapping spaces in
  quasi-categories. \emph{Algebr. Geom. Topol.}, \textbf{11}, 263--325.

\bibitem{bib37} Dwyer, W. G., and Kan, D. M. 1980. Simplicial localizations of
  categories. \emph{J. Pure Appl. Algebra}, \textbf{17}, 267--284.

\bibitem{bib38} Dwyer, W. G., Kan, D. M., and Smith, J. H. 1989. Homotopy
  commutative diagrams and their realizations. \emph{J. Pure Appl. Algebra},
  \textbf{57}, 5--24.

\bibitem{bib39} Dwyer, W. G., Hirschhorn, P. S., Kan, D. M., and Smith, J. H.
  2004. \emph{Homotopy Limit Functors on Model Categories and Homotopical
  Categories}. Mathematical Surveys and Monographs, vol.~113. American
  Mathematical Society.

\bibitem{bib40} Ehlers, P. J., and Porter, T. 2000. Joins for (augmented)
  simplicial sets. \emph{J. Pure Appl. Algebra}, \textbf{145}, 37--44.

\bibitem{bib41} Eilenberg, S., and Kelly, G. M. 1966. Closed categories. Page
  421--562 of: \emph{Proceedings of the Conference on Categorical Algebra}
  (La Jolla 1965). Springer-Verlag.

\bibitem{bib42} Eilenberg, S., and Mac Lane, S. 1945. General theory of natural
  equivalences. \emph{Trans. Amer. Math. Soc.}, \textbf{58}, 231--294.

\bibitem{bib43} Freyd, P. J. 1976. Properties invariant within equivalence types
  of categories. Page 55--61 of: Heller, A., and Tierney, M. (eds),
  \emph{Algebra, Topology and Category Theories}. Academic Press.

\bibitem{bib44} Gabriel, P., and Zisman, M. 1967. \emph{Calculus of Fractions and
  Homotopy Theory}. Ergebnisse der Math., no.~35. Springer-Verlag.

\bibitem{bib45} Gambino, N. 2010. Weighted limits in simplicial homotopy theory.
  \emph{J. Pure Appl. Algebra}, \textbf{214}(7), 1193--1199.

\bibitem{bib46} Gepner, D., Haugseng, R., and Kock, J. 2021. $\infty$-operads as
  analytic monads. \emph{International Mathematics Research Notices}.

\bibitem{bib47} Grandis, M., and Paré, R. 2004. Adjoint for double categories.
  \emph{Cah. Topol. Géom. Différ. Catég.}, \textbf{45}(3), 193--240.

\bibitem{bib48} Gray, J. W. 1966. Fibred and cofibred categories. Page 21--83
  of: Eilenberg, S., Harrison, D. K., Röhrl, H., and MacLane, S. (eds),
  \emph{Proceedings of the Conference on Categorical Algebra} (La Jolla 1965).
  Springer-Verlag.

\bibitem{bib49} Groth, M. 2016. Characterizations of abstract stable homotopy
  theories. arXiv:1602.07632.

\bibitem{bib50} Guillou, B., and May, J. P. 2011. Enriched model categories and
  presheaf categories. arXiv:1110.3567.

\bibitem{bib51} Harpaz, Y. 2017. Introduction to stable $\infty$-categories.
  www.math.univ-paris13.fr/\textasciitilde{}harpaz/stable\_infinity\_categories.pdf.

\bibitem{bib52} Henry, S. 2020. The language of a model category.
  www.uwo.ca/math/faculty/kapulkin/seminars/hottfiles/Henry-2020-01-23-HoTTEST.pdf.

\bibitem{bib53} Hermida, C. 2000. Representable multicategories. \emph{Adv.
  Math.}, \textbf{151}, 164--225.

\bibitem{bib54} Hess, K., Kędziorek, M., Riehl, E., and Shipley, B. 2017. A
  necessary and sufficient condition for induced model structures. \emph{J.
  Topol.}, \textbf{10}(2), 324--369.

\bibitem{bib55} Hirschhorn, P. 2003. \emph{Model Categories and Their
  Localizations}. Mathematical Surveys and Monographs, no.~99. American
  Mathematical Society.

\bibitem{bib56} Hirschowitz, A., and Simpson, C. 1998. Descente pour les
  $n$-champs. arXiv:math/9807049.

\bibitem{bib57} Hovey, M. 1999. \emph{Model Categories}. Mathematical Surveys and
  Monographs, vol.~63. American Mathematical Society.

\bibitem{bib58} Johnson, N., and Yau, D. 2019. A bicategorical pasting theorem.
  arXiv:1910.01220.

\bibitem{bib59} Johnson, N., and Yau, D. 2021. \emph{2-Dimensional Categories}.
  Oxford University Press.

\bibitem{bib60} Joyal, A. 1997. Disks, duality and $\Theta$-categories.
  ncatlab.org/nlab/files/JoyalTheta-Categories.pdf.

\bibitem{bib61} Joyal, A. 2002. Quasi-categories and Kan complexes. \emph{J.
  Pure Appl. Algebra}, \textbf{175}, 207--222.

\bibitem{bib62} Joyal, A. 2008a. Notes on quasi-categories.
  www.math.uchicago.edu/\textasciitilde{}may/IMA/Joyal.pdf.

\bibitem{bib63} Joyal, A. 2008b. \emph{The Theory of Quasi-Categories and its
  Applications}. Quadern 45 vol II. Centre de Recerca Matemàtica Barcelona.

\bibitem{bib64} Joyal, A., and Tierney, M. 2007. Quasi-categories vs Segal
  spaces. Page 277--325 of: Davydov, A. (ed), \emph{Categories in Algebra,
  Geometry and Mathematical Physics (StreetFest)}. Contemp. Math., vol.~431.
  American Mathematical Society.

\bibitem{bib65} Kazhdan, D., and Varshavsky, Y. 2014. Yoneda lemma for complete
  Segal spaces. \emph{Funct Anal Its Appl}, \textbf{48}(2), 81--106.

\bibitem{bib66} Kelly, G.M. 1964. On Mac Lane's conditions for coherence of
  natural associativities, commutativities, etc. \emph{J. Algebra},
  \textbf{1}, 397--402.

\bibitem{bib67} Kelly, G. M. 1989. Elementary observations on 2-categorical
  limits. \emph{Bull. Aust. Math. Soc.}, \textbf{49}, 301--317.

\bibitem{bib68} Kelly, G. M. 2005. Basic concepts of enriched category theory.
  \emph{Repr. Theory Appl. Categ.}, \textbf{10}.

\bibitem{bib69} Kelly, G. M., and Mac Lane, S. 1971. Coherence in Closed
  Categories. \emph{J. Pure Appl. Algebra}, \textbf{1}(1), 97--140.

\bibitem{bib70} Kelly, G. M., and Street, R. H. 1974. \emph{Review of the
  Elements of 2-Categories}. Lecture Notes in Mathematics, vol.~420.
  Springer-Verlag. Page 75--103.

\bibitem{bib71} Lack, S. 2002. A Quillen model structure for 2-categories.
  \emph{J. K-Theory}, \textbf{26}(2), 171--205.

\bibitem{bib72} Lack, S. 2007. Homotopy-theoretic aspects of 2-monads. \emph{J.
  Homotopy Relat. Struct.}, \textbf{2}(2), 229--260.

\bibitem{bib73} Lack, Stephen. 2010. A 2-Categories Companion. In: May, J. P.,
  and Baez, J. C. (eds), \emph{Towards Higher Categories}. The IMA Volumes in
  Mathematics and its Applications, vol.~152. Springer.

\bibitem{bib74} Lawvere, F. W. 2005. An elementary theory of the category of sets
  (long version) with commentary. \emph{Repr. Theory Appl. Categ.},
  \textbf{11}, 1--35.

\bibitem{bib75} Leinster, T. 2002. Generalized enrichment of categories.
  \emph{J. Pure Appl. Algebra}, \textbf{168}, 391--406.

\bibitem{bib76} Leinster, T. 2004. \emph{Higher Operads, Higher Categories}.
  London Mathematical Society Lecture Note Series, vol.~298. Cambridge
  University Press.

\bibitem{bib77} Lurie, J. 2006 (May). \emph{Stable Infinity Categories}.
  arXiv:math/0608228.

\bibitem{bib78} Lurie, J. 2009a. \emph{Higher Topos Theory}. Annals of
  Mathematical Studies, vol.~170. Princeton University Press.

\bibitem{bib79} Lurie, J. 2009b. \emph{($\infty$,2)-Categories and the Goodwillie
  Calculus I}. arXiv:0905.0462.

\bibitem{bib80} Lurie, J. 2017 (September). \emph{Higher Algebra}.
  www.math.ias.edu/\textasciitilde{}lurie/papers/HA.pdf.

\bibitem{bib81} Mac Lane, S. 1998. \emph{Categories for the Working
  Mathematician}. Second edn. Graduate Texts in Mathematics. Springer-Verlag.

\bibitem{bib82} Makkai, M. 1995. \emph{First Order Logic with Dependent Sorts,
  with Applications to Category Theory}. www.math.mcgill.ca/makkai/.

\bibitem{bib83} Makkai, M., and Paré, R. 1989. \emph{Accessible Categories: The
  Foundations of Categorical Model Theory}. Contemp. Math., vol.~104.
  American Mathematical Society.

\bibitem{bib84} May, J. P., and Ponto, K. 2012. \emph{More concise algebraic
  topology: Localization, Completion, and Model Categories}. Chicago Lectures
  in Mathematics. University of Chicago Press.

\bibitem{bib85} Mazel-Gee, A. 2016. Quillen adjunctions induce adjunctions of
  quasicategories. \emph{New York J Math.}, \textbf{22}, 57--93.

\bibitem{bib86} Mazel-Gee, A. 2019. A user's guide to co/cartesian fibrations.
  \emph{Graduate Journal of Mathematics}, \textbf{4}(1), 42--53.

\bibitem{bib87} Myers, D. J. 2018. \emph{String Diagrams For Double Categories
  and Equipments}. arXiv:1612.02762.

\bibitem{bib88} Nguyen, H. K., Raptis, G., and Schrade, C. 2020. Adjoint functor
  theorems for $\infty$-categories. \emph{J. Lond. Math. Soc.},
  \textbf{101}(2), 659--681.

\bibitem{bib89} Ozornova, V., and Rovelli, M. 2020. Model structures for
  ($\infty,n$)-categories on (pre)stratified simplicial sets and prestratified
  simplicial spaces. \emph{Algebr. Geom. Topol.}, \textbf{20}, 1543--1600.

\bibitem{bib90} Pellissier, R. 2002. \emph{Catégories enrichies faibles}. Ph.D.
  thesis, Université de Nice-Sophia Antipolis.

\bibitem{bib91} Power, A. J. 1990. A 2-categorical pasting theorem. \emph{J.
  Algebra}, \textbf{129}, 439--445.

\bibitem{bib92} Prller, A. 1985. A language for category theory in which natural
  equivalence implies elementary equivalence of models. \emph{Zeitschrift f.
  math. Logik und Grundlagen Math.}, \textbf{31}, 227--234.

\bibitem{bib93} Quillen, D. G. 1967. \emph{Homotopical Algebra}. Lecture Notes in
  Mathematics, vol.~43. Springer-Verlag.

\bibitem{bib94} Radulescu-Banu, A. 2009. Cofibrations in homotopy theory.
  arXiv:math/0610009.

\bibitem{bib95} Rasekh, N. 2017a. Cartesian fibrations and representability.
  arXiv:1711.03670.

\bibitem{bib96} Rasekh, N. 2017b. Yoneda lemma for simplicial spaces.
  arXiv:1711.03160.

\bibitem{bib97} Rasekh, N. 2021a. Quasi-categories vs. Segal spaces: cartesian
  edition. \emph{J. Homotopy Relat. Struct.}.

\bibitem{bib98} Rasekh, Nima. 2021b. Cartesian fibrations of complete Segal
  spaces. arXiv:2102.05190.

\bibitem{bib99} Reedy, C. 1974. Homotopy theory of model categories.
  www-math.mit.edu/\textasciitilde{}psh/reedy.pdf.

\bibitem{bib100} Rezk, C. 2001. A model for the homotopy theory of homotopy
  theory. \emph{Trans. Amer. Math. Soc.}, \textbf{353}(3), 973--1007.

\bibitem{bib101} Rezk, C. 2010. A cartesian presentation of weak $n$-categories.
  \emph{Geom. Topol.}, \textbf{14}(1), 521--571.

\bibitem{bib102} Rezk, C. 2019. \emph{Stuff about quasi-categories}.
  faculty.math.illinois.edu/\textasciitilde{}rezk/quasicats.pdf.

\bibitem{bib103} Riehl, E. 2014. \emph{Categorical Homotopy Theory}. New
  Mathematical Monographs, vol.~24. Cambridge University Press.

\bibitem{bib104} Riehl, E. 2016. \emph{Category Theory in Context}. Aurora
  Modern Math Originals. Dover Publications.

\bibitem{bib105} Riehl, E. 2017. Complicial sets, an overture. In: 2016
  \emph{MATRIX Annals}. MATRIX Book Series, no.~1--28. Springer.

\bibitem{bib106} Riehl, E., and Shulman, M. 2017. A type theory for synthetic
  $\infty$-categories. \emph{High. Struct.}, \textbf{1}(1), 147--224.

\bibitem{bib107} Riehl, E., and Verity, D. 2014. The theory and practice of
  Reedy categories. \emph{Theory Appl. Categ.}, \textbf{29}(9), 256--301.

\bibitem{bib108} Riehl, E., and Verity, D. 2015. The 2-category theory of
  quasi-categories. \emph{Adv. Math.}, \textbf{280}, 549--642.

\bibitem{bib109} Riehl, E., and Verity, D. 2016. Homotopy coherent adjunctions
  and the formal theory of monads. \emph{Adv. Math.}, \textbf{286}, 802--888.

\bibitem{bib110} Riehl, E., and Verity, D. 2017. Fibrations and Yoneda's lemma
  in an $\infty$-cosmos. \emph{J. Pure Appl. Algebra}, \textbf{221}(3),
  499--564.

\bibitem{bib111} Riehl, E., and Verity, D. 2018. The comprehension construction.
  \emph{High. Struct.}, \textbf{2}(1), 116--190.

\bibitem{bib112} Riehl, E., and Verity, D. 2020a. On the construction of limits
  and colimits in $\infty$-categories. \emph{Theory Appl. Categ.},
  \textbf{35}(30), 1101--1158. arXiv:1808.09835.

\bibitem{bib113} Riehl, E., and Verity, D. 2020b. Recognizing quasi-categorical
  limits and colimits in homotopy coherent nerves. \emph{Appl. Categ.
  Struct.}, \textbf{28}(4), 669--716. arXiv:1808.09834.

\bibitem{bib114} Riehl, E., and Verity, D. 2021. Cartesian exponentiation and
  monadicity. arXiv:2101.09853.

\bibitem{bib115} Shah, Jay. 2018. Parametrized higher category theory and higher
  algebra: Exposé II - Indexed homotopy limits and colimits. arXiv:1809.05892.

\bibitem{bib116} Simpson, C. 2012. \emph{Homotopy Theory of Higher Categories}.
  New Mathematical Monographs, no.~19. Cambridge University Press.

\bibitem{bib117} Street, R. H. 1974a. Elementary cosmoi I. Page 134--180 of:
  \emph{Category Seminar} (Proc. Sem., Sydney, 1972/1973). Lecture Notes in
  Math., vol.~420. Springer.

\bibitem{bib118} Street, R. H. 1974b. Fibrations and Yoneda's lemma in a
  2-category. Page 104--133, of: \emph{Category Seminar} (Proc. Sem., Sydney,
  1972/1973). Lecture Notes in Math., vol.~420. Springer.

\bibitem{bib119} Street, R. H. 1980. Fibrations in Bicategories. \emph{Cah.
  Topol. Géom. Différ. Catég.}, \textbf{XXI-2}, 111--159.

\bibitem{bib120} Street, R. H. 1987a. The algebra of oriented simplexes.
  \emph{J. Pure Appl. Algebra}, \textbf{49}, 283--335.

\bibitem{bib121} Street, R. H. 1987b. Correction to ``Fibrations in
  bicategories''. \emph{Cah. Topol. Géom. Différ. Catég.}, \textbf{28}(1),
  53--56.

\bibitem{bib122} Street, R. H. 2003. Weak omega-categories. Page 207--213 of:
  \emph{Diagrammatic Morphisms and Applications}. Contemp. Math., vol.~318.
  American Mathematical Society.

\bibitem{bib123} Street, R., and Walters, R. 1978. Yoneda structures on
  2-categories. \emph{J. Algebra}, \textbf{50}, 350--379.

\bibitem{bib124} Sulyma, Y. J. F. 2017. $\infty$-categorical monadicity and
  descent. \emph{New York Journal of Mathematics}, \textbf{23}, 749--777.

\bibitem{bib125} The Univalent Foundations Program. 2013. \emph{Homotopy Type
  Theory: Univalent Foundations of Mathematics}. Institute for Advanced Study:
  homotopytypetheory.org/book.

\bibitem{bib126} Ulmer, F. 1968. Properties of dense and relative adjoint
  functors. \emph{J. Algebra}, \textbf{3}, 77--95.

\bibitem{bib127} Verdier, J. L. 1996. \emph{Des Catégories Dérivées des
  Catégories Abéliennes}. \emph{Astérisque}, \textbf{239}, 1--253+ix.

\bibitem{bib128} Verity, D. 2008a. \emph{Complicial Sets, Characterising the
  Simplicial Nerves of Strict $\omega$-Categories}. \emph{Mem. Amer. Math.
  Soc.}, \textbf{193}(905).

\bibitem{bib129} Verity, D. 2008b. Weak complicial sets I, basic homotopy
  theory. \emph{Adv. Math.}, \textbf{219}, 1081--1149.

\bibitem{bib130} Verity, D. 2011. \emph{Enriched Categories, Internal Categories
  and Change of Base}. \emph{Repr. Theory Appl. Categ.}, \textbf{20}, 1--266.

\bibitem{bib131} Weber, M. 2007. Yoneda structures from 2-toposes. \emph{Appl.
  Categ. Structures}, \textbf{15}(3), 259--323.

\bibitem{bib132} Wood, R. J. 1982. Abstract proarrows I. \emph{Cah. Topol.
  Géom. Différ. Catég.}, \textbf{XXIII-3}, 279--290.

\bibitem{bib133} Zaganidis, D. 2017. \emph{Towards an ($\infty$,2)-category of
  Homotopy Coherent Monads in an $\infty$-cosmos}. Ph.D. thesis, École
  Polytechnique Fédérale de Lausanne.

\end{thebibliography}
